\documentclass{article}
\usepackage[utf8]{inputenc}
\usepackage[english]{babel}
\usepackage{geometry}
\usepackage{amsmath}
\usepackage{graphicx}
\usepackage{hyperref}
\usepackage{array}
\usepackage{algorithmic}
\usepackage{listings}

\geometry{top=2.5cm, bottom=2.5cm, left=2.5cm, right=2.5cm}

\begin{document}
\begin{titlepage}
    \begin{center}
        \vspace*{2cm}

        \Huge
        \textbf{Advanced Databases}
        \vspace{0.25cm}

        \LARGE
        \textbf{INFO-H-415}

        \vspace{0.25cm}
        \LARGE
        Object-oriented databases with ObjectDB and ObjectBox

        \vspace{0.25cm}
        \textit{14 December 2026}

        \vspace{2cm}
           \Large
        \textbf{Amine Ayadi}\\
		\textbf{Martin Gouverneur}\\
		\textbf{Filipe Pereira Martins}\\
		\textbf{Ethan Van Ruyskensvelde}\\

        \vspace{2cm}

        \textsc{Université Libre de Bruxelles (ULB)}

    \end{center}
\end{titlepage}

\tableofcontents
\newpage

\section{Introduction}
The goal of this project is to study object-oriented databases and compare two database management systems, ObjectDB and ObjectBox. Object-oriented databases are designed to store and manipulate objects directly, without needing to map application objects to relational tables. This project will present the fundamental concepts of object-oriented database technology and examine how they are implemented in ObjectDB and ObjectBox. A common application will be developed using both systems to perform a set of representative queries and updates. The performance of these operations will then be evaluated on datasets of different sizes to analyze the scalability of each system. Finally, the results will be compared with those obtained using a traditional relational database management system, in order to see the advantages and limitations of object-oriented databases, and determine if it is the technology of choice for the application.

\section{Foundations of Object-Oriented Databases}
\subsection{Origins and motivation}
\subsection{Fundamental Principles of Object-Oriented Database Management Systems}
\subsection{Standards and norms}
\subsection{Relational Model (RDBMS) vs. Object-Oriented Model (OODBMS)}

\section{Overview of the DBMS Studied}
\subsection{ObjectDB: Object DBMS for Java / JPA}
\subsection{ObjectBox: Embedded Object-Oriented and High-Performance NoSQL DBMS}
\subsection{PostgreSQL: The benchmark relational DBMS (Baseline)}

\section{Use Cases and Domain Modeling}
\subsection{Description of the application domain}
\subsection{Conceptual schema and data models}
\subsection{Software architecture}
\subsection{Dataset generation}

\section{Experimentation and Benchmarking Protocol}
\subsection{Benchmark environment}
\subsection{Measurement methodology}
\subsection{Evaluated operations}

\section{Results and Comparative Analysis}
\subsection{Summary of results}
\subsection{Scalability analysis}
\subsection{Critical analysis by DBMS}
\subsubsection{PostgreSQL}
\subsubsection{ObjectDB}
\subsubsection{ObjectBox}

\section{Overall Assessment: Pros and Cons}
\subsection{Comparison table}
\subsection{Usage recommendations}

\section{Conclusion}

\section{Bibliography and Appendices}

\end{document}
